% !TEX root = ../main.tex
\questiontitle{Question 3: Forest-Fire Monitoring Wireless Sensor Network}

\begin{question}{3}
A forest-fire monitoring system consists of 1,000 battery-powered sensor nodes randomly deployed over a large
forest. Each node measures temperature, humidity and smoke concentration. The nodes forward information to a base
station located at the edge of the forest.
\end{question}

% ------------------------------------------------------------------ (a)
\begin{question}{3 (a)}
Would you select a star/single-hop or multi-hop mesh topology? Justify your choice from the perspective of energy
consumption and coverage.
\end{question}

\answerhead{(a) Star/single-hop versus multi-hop mesh topology}
A multi-hop mesh topology is more suitable for the forest-fire monitoring system.

\begin{itemize}
\item \textbf{Energy consumption:} Distant sensors can forward data through nearby nodes instead of transmitting
directly over long distances, potentially reducing transmission energy.
\item \textbf{Coverage:} Multi-hop communication extends connectivity across a large forest, including locations
outside the base station's direct communication range.
\item \textbf{Reliability:} Multiple possible routes can provide alternative paths when individual nodes or links
fail.
\end{itemize}

A single-hop star topology would require every sensor to reach the base station directly, potentially consuming
more energy and limiting coverage.

% ------------------------------------------------------------------ (b)
\begin{question}{3 (b)}
Would a flat or hierarchical/clustered WSN (Wireless Sensor Network) architecture be more suitable? Explain your
choice.
\end{question}

\answerhead{(b) Flat versus hierarchical/clustered WSN (Wireless Sensor Network) architecture}
A hierarchical/clustered architecture is more suitable for 1,000 randomly deployed sensors.

The network is divided into clusters, each managed by a cluster head. Ordinary sensor nodes send their readings to
their cluster head, which aggregates the data and forwards it toward the base station.

Advantages include:
\begin{itemize}
\item Reduced redundant transmissions.
\item Better scalability for a large sensor network.
\item Lower communication overhead for ordinary nodes.
\item More efficient coordination and energy management.
\end{itemize}

Cluster heads should be rotated or selected using residual energy and communication cost to avoid premature
battery depletion.

\simfigure[width=0.72\linewidth]{../media/figures/q3_architecture.pdf}{Hierarchical/clustered multi-hop WSN (Wireless
Sensor Network) for the forest: 1,000 sensor nodes over a 2 kilometre $\times$ 2 kilometre forest. Each cluster head
aggregates its members' readings and forwards them to the base station over a multi-hop route.}

% ------------------------------------------------------------------ (c)
\begin{question}{3 (c)}
Explain how data aggregation at intermediate nodes or cluster heads can extend network lifetime.
\end{question}

\answerhead{(c) How data aggregation extends network lifetime}
Data aggregation combines readings from multiple sensors into a compact summary before forwarding them. For
example, a cluster head can calculate average temperature and humidity while preserving maximum temperature and
smoke readings that may indicate a fire.

Without aggregation, many individual packets must travel toward the base station. With aggregation, fewer packets
are transmitted over long distances, reducing communication energy and congestion. This helps conserve battery
power and extends network lifetime while retaining important fire-detection information.
